\begin{center}ALLGEMEINE FORMELN FÜR DIE WIRKUNG EINES LEUCHTENDEN PUNKTS $P$ AUF EINEN PUNKT $p$.\end{center}
I. Der Punkt $P$ sei von $p$ durch eine geschlossene oder unendliche Fläche getrennt. Ihr offener Teil heiße $s$. Es sei $\mathrm ds$ ein Element von $s$ und $R,r$ seine Abstände von $P,p$. Es sei $\rho$ eine unbestimmte Normale zu $\mathrm ds$, die zu der Seite weist, auf der $p$ liegt, $\lambda$ eine Wellenlänge, $\dfrac{2\pi i}{\lambda}=\alpha$ und $w$ der Winkel zwischen $R$ und $r$. Dann wird der Schwingungszustand in $p$ durch
\[
\int\frac{\mathrm ds}{Rr}\left(\frac{\partial R}{\partial\rho}+\frac{\partial r}{\partial\rho}\right)\frac{e^{\alpha(R+r)}}{\sin w}
\]
ausgedrückt, wobei sich das Integral über alle Teile von $s$ erstreckt. Offenbar sind hier $\dfrac{\partial R}{\partial\rho},-\dfrac{\partial r}{\partial\rho}$ die Sinus der Neigungswinkel von $R$ und $r$ gegenüber $\mathrm ds$.
\SourceLeaf{642}

II. Der Flächenbereich $s$ sei durch die Linie $u$ begrenzt. Es sei $\mathrm du$ ein. Element von $u$ und $R,r$ seine Abstände von $P$ und $p$. Es sei $w$ der Winkel zwischen $R$ und $r$ und $v$ der Winkel zwischen $u$ und der Ebene durch $R,r$.

Dann ist der Schwingungszustand in $p$
\[
=\int\frac{\mathrm du.\sin v}{Rr.\sin w}e^{\alpha(R+r)}
\]
Kürzer gefasst

Es seien $x,y,z$ die Koordinaten jedes Punkts im Raum. Für die beiden Punkte gelte
\[
\begin{array}{lccc}
P\, .\, .\, .&0,&0,&0\\
p\, .\, .\, .&0,&0,&h
\end{array}
\]
Der Schwingungszustand in $p$ ist dann
\[
=\int\frac{x\mathrm dy-y\mathrm dx}{xx+yy}.e^{\alpha(R+r)}:h
\]
wobei sich das Integral über die Randlinie erstreckt. Oder kürzer, wenn $\dfrac yx=\operatorname{tg}\theta$,
\[
=\int e^{\alpha(R+r)}\mathrm d\theta:h
\]

